\documentclass[11pt]{article}
\usepackage[utf8]{inputenc}
\usepackage[T1]{fontenc}
\usepackage{amsmath}
\usepackage{amssymb}
\usepackage{multicol}
\usepackage{geometry}
\usepackage{tikz}
\usepackage{enumitem}
\usepackage{xcolor}
\usepackage{titlesec}

% Configurações de layout
\geometry{a4paper, left=1cm, right=1cm, top=0.5cm, bottom=1.2cm}
\setlength{\columnseprule}{0.4pt}

% Cores para títulos
\titleformat{\section}{\normalfont\Large\bfseries\color{blue}}{\thesection}{1em}{}
\titleformat{\subsection}{\normalfont\large\bfseries\color{red}}{\thesubsection}{1em}{}

% Reduz espaço no multicols
\setlength{\multicolsep}{0pt}

\title{\textcolor{green}{Revisão: Medidas de Dispersão \\
    Amplitude, Variância e Desvio Padrão}}
\author{Professor: Jefferson}
\date{}

\begin{document}

\maketitle
\vspace{-0.8cm}

\begin{multicols}{2}

\section*{Instruções}
\textbf{Copie e responda no caderno.}


\begin{enumerate}[leftmargin=*]

\item Calcule a amplitude dos conjuntos:
\begin{enumerate}[label=\alph*)]
\item 5, 8, 12, 15, 20
\item 10, 20, 30, 40, 50
\item 3, 7, 11, 15, 19, 23
\item 100, 150, 200, 250, 300
\end{enumerate}

\item Calcule a média e a amplitude:
\begin{enumerate}[label=\alph*)]
\item 6, 9, 12, 15, 18
\item 2, 5, 8, 11, 14, 17
\end{enumerate}

\item Calcule a variância e o desvio padrão do conjunto: 4, 6, 8, 10, 12

\item Calcule a variância e o desvio padrão do conjunto: 2, 4, 6, 8, 10

\item Compare o desvio padrão dos dois conjuntos abaixo. Qual é mais disperso?
\begin{enumerate}[label=\alph*)]
\item Conjunto A: 5, 10, 15, 20, 25
\item Conjunto B: 12, 13, 14, 15, 16
\end{enumerate}

\item Dado o conjunto: 10, 12, 14, 16, 18, 20
\begin{enumerate}[label=\alph*)]
\item Calcule a amplitude
\item Calcule a média
\item Calcule a variância
\item Calcule o desvio padrão
\end{enumerate}

\end{enumerate}

\begin{enumerate}[leftmargin=*, start=7]

\item As notas de dois alunos foram:
\begin{itemize}
\item Aluno A: 5, 6, 7, 8, 9
\item Aluno B: 2, 5, 7, 9, 12
\end{itemize}
\begin{enumerate}[label=\alph*)]
\item Calcule a média de cada um
\item Calcule o desvio padrão de cada um
\item Qual aluno teve melhor desempenho? Qual foi mais regular?
\end{enumerate}

\item Complete a tabela:

\begin{center}
% Centralizar a tabela a esquerda com 
\hspace*{-0.9cm}
% Diminuir a tabela para caber melhor
% opções: \small, \footnotesize, \scriptsize
\footnotesize
\setlength{\tabcolsep}{4pt}
\renewcommand{\arraystretch}{0.8}
\begin{tabular}{|c|c|c|c|c|}
\hline
\textbf{Dados} & \textbf{Amplitude} & \textbf{Média} & \textbf{Variância} & \textbf{Desvio Padrão} \\
\hline
2, 4, 6, 8, 10 & & & & \\
\hline
1, 3, 5, 7, 9 & & & & \\
\hline
10, 20, 30, 40, 50 & & & & \\
\hline
\end{tabular}
\end{center}

\item As temperaturas de duas cidades foram registradas:

\textbf{Cidade X:} 18, 20, 22, 24, 26 (em °C)

\textbf{Cidade Y:} 15, 20, 22, 25, 28

\begin{enumerate}[label=\alph*)]
\item Calcule a média de cada cidade
\item Calcule o desvio padrão de cada cidade
\item Qual cidade tem clima mais regular? Por quê?
\end{enumerate}

\item Em uma fábrica, o controle de qualidade mediu o diâmetro (em mm) de 6 peças: 48, 50, 52, 49, 51, 50
\begin{enumerate}[label=\alph*)]
\item Calcule a média
\item Calcule a amplitude
\item Calcule o desvio padrão
\end{enumerate}

\item Uma prova foi aplicada em duas turmas:

\textbf{Turma A:} notas: 4, 5, 6, 7, 8, 9, 10

\textbf{Turma B:} notas: 5, 6, 7, 7, 7, 8, 9

\begin{enumerate}[label=\alph*)]
\item Calcule a média de cada turma
\item Calcule o desvio padrão de cada turma
\item Qual turma teve notas mais homogêneas?
\end{enumerate}

\item Dados os conjuntos:

$X = \{10, 20, 30, 40, 50\}$ e $Y = \{28, 29, 30, 31, 32\}$

\begin{enumerate}[label=\alph*)]
\item Calcule a média de X e Y
\item Calcule o desvio padrão de X e Y
\item O que você observa comparando os resultados?
\end{enumerate}

\end{enumerate}

\end{multicols}

\end{document}
